% Lecture 10 - Page 6 (6-XII-1994).
% Fragment for % Lecture 10 - Page 6 (6-XII-1994).
% Fragment for % Lecture 10 - Page 6 (6-XII-1994).
% Fragment for % Lecture 10 - Page 6 (6-XII-1994).
% Fragment for \input{courses/pages/curs10-pagina6-en}.
% No preamble and no forced page break.

\subsection*{Lecture 10 - Page 6}

\noindent\underline{Def.}: A group $G$ is called a \textbf{$p$-group},
$p$ prime, if:
\[
(\forall)\,a\in G,\qquad o(a)=p^t,\quad t\geq 0.
\]
% In the manuscript, the term "p-group" in the opening sentence is crossed out,
% although it is the term defined by the condition that follows.

\medskip
\noindent\underline{Remark}: If $G$ is a finite group, then
\[
\boxed{G\text{ is a }p\text{-group}\quad\Longleftrightarrow\quad
\operatorname{ord}(G)=p^n}.
\]

\noindent\underline{Proof}: ``$\Leftarrow$'' evident.

\noindent ``$\Rightarrow$'': By contradiction, suppose
$\operatorname{ord}(G)\neq p^n$.
Then there is a prime $q\neq p$ such that $q\mid\operatorname{ord}(G)$.

\noindent By Cauchy's theorem, $(\exists)\,a\in G$ with $o(a)=q$;
since $G$ is a $p$-group, this gives $p=q$, a contradiction.

\bigskip
\noindent-- Let
\[
U_{p^n}=\{z\in\mathbb{C}\mid z^{p^n}=1\},
\qquad U_{p^n}\leq(\mathbb{C}^{*},\cdot).
\]
\noindent Also $\operatorname{ord}(U_{p^n})=p^n$.

\[
\bigcup_{n\geq 1}U_{p^n}\ \leq\ (\mathbb{C}^{*},\cdot).
\]
\noindent This is an infinite $p$-group.

\bigskip
\noindent\underline{Def.}: Let $G$ be a group and $H\leq G$.
$H$ is a \underline{$p$-subgroup} of $G$ if and only if
$(H,\cdot)$ is a $p$-group.

\bigskip
\noindent\fbox{\textbf{Sylow's First Theorem}}.
Let $G$ be a finite group with $\operatorname{ord}(G)=p^t m$,
$t\geq 1$ and $p\nmid m$.

\medskip
\noindent Then for every $i=\overline{1,t}$ there exists $H\leq G$
with $|H|=p^i$.

\medskip
\noindent Moreover, if $1\leq i<t$, there exists $H\leq G$,
$|H|=p^i$, which is normal in a subgroup of $G$
of order $p^{i+1}$.

\bigskip
\noindent\underline{Proof}: For $i=1$, O.K. by Cauchy's theorem.

\medskip
\noindent Under the same hypotheses, a $p$-subgroup of $G$
of order $p^t$ is called a \underline{Sylow $p$-subgroup} of $G$.
\quad (Notation: $A=\operatorname{Syl}_p(G)$.)

\bigskip
\noindent\underline{Remark}: On $A=\operatorname{Syl}_p(G)$,
$G\times A\longrightarrow A$ acts by conjugation:
\[
P\in A,\ a\in G,\qquad
(a,P)\longmapsto aPa^{-1}
\quad\text{(the conjugate of $P$ by $a$)}.
\]
.
% No preamble and no forced page break.

\subsection*{Lecture 10 - Page 6}

\noindent\underline{Def.}: A group $G$ is called a \textbf{$p$-group},
$p$ prime, if:
\[
(\forall)\,a\in G,\qquad o(a)=p^t,\quad t\geq 0.
\]
% In the manuscript, the term "p-group" in the opening sentence is crossed out,
% although it is the term defined by the condition that follows.

\medskip
\noindent\underline{Remark}: If $G$ is a finite group, then
\[
\boxed{G\text{ is a }p\text{-group}\quad\Longleftrightarrow\quad
\operatorname{ord}(G)=p^n}.
\]

\noindent\underline{Proof}: ``$\Leftarrow$'' evident.

\noindent ``$\Rightarrow$'': By contradiction, suppose
$\operatorname{ord}(G)\neq p^n$.
Then there is a prime $q\neq p$ such that $q\mid\operatorname{ord}(G)$.

\noindent By Cauchy's theorem, $(\exists)\,a\in G$ with $o(a)=q$;
since $G$ is a $p$-group, this gives $p=q$, a contradiction.

\bigskip
\noindent-- Let
\[
U_{p^n}=\{z\in\mathbb{C}\mid z^{p^n}=1\},
\qquad U_{p^n}\leq(\mathbb{C}^{*},\cdot).
\]
\noindent Also $\operatorname{ord}(U_{p^n})=p^n$.

\[
\bigcup_{n\geq 1}U_{p^n}\ \leq\ (\mathbb{C}^{*},\cdot).
\]
\noindent This is an infinite $p$-group.

\bigskip
\noindent\underline{Def.}: Let $G$ be a group and $H\leq G$.
$H$ is a \underline{$p$-subgroup} of $G$ if and only if
$(H,\cdot)$ is a $p$-group.

\bigskip
\noindent\fbox{\textbf{Sylow's First Theorem}}.
Let $G$ be a finite group with $\operatorname{ord}(G)=p^t m$,
$t\geq 1$ and $p\nmid m$.

\medskip
\noindent Then for every $i=\overline{1,t}$ there exists $H\leq G$
with $|H|=p^i$.

\medskip
\noindent Moreover, if $1\leq i<t$, there exists $H\leq G$,
$|H|=p^i$, which is normal in a subgroup of $G$
of order $p^{i+1}$.

\bigskip
\noindent\underline{Proof}: For $i=1$, O.K. by Cauchy's theorem.

\medskip
\noindent Under the same hypotheses, a $p$-subgroup of $G$
of order $p^t$ is called a \underline{Sylow $p$-subgroup} of $G$.
\quad (Notation: $A=\operatorname{Syl}_p(G)$.)

\bigskip
\noindent\underline{Remark}: On $A=\operatorname{Syl}_p(G)$,
$G\times A\longrightarrow A$ acts by conjugation:
\[
P\in A,\ a\in G,\qquad
(a,P)\longmapsto aPa^{-1}
\quad\text{(the conjugate of $P$ by $a$)}.
\]
.
% No preamble and no forced page break.

\subsection*{Lecture 10 - Page 6}

\noindent\underline{Def.}: A group $G$ is called a \textbf{$p$-group},
$p$ prime, if:
\[
(\forall)\,a\in G,\qquad o(a)=p^t,\quad t\geq 0.
\]
% In the manuscript, the term "p-group" in the opening sentence is crossed out,
% although it is the term defined by the condition that follows.

\medskip
\noindent\underline{Remark}: If $G$ is a finite group, then
\[
\boxed{G\text{ is a }p\text{-group}\quad\Longleftrightarrow\quad
\operatorname{ord}(G)=p^n}.
\]

\noindent\underline{Proof}: ``$\Leftarrow$'' evident.

\noindent ``$\Rightarrow$'': By contradiction, suppose
$\operatorname{ord}(G)\neq p^n$.
Then there is a prime $q\neq p$ such that $q\mid\operatorname{ord}(G)$.

\noindent By Cauchy's theorem, $(\exists)\,a\in G$ with $o(a)=q$;
since $G$ is a $p$-group, this gives $p=q$, a contradiction.

\bigskip
\noindent-- Let
\[
U_{p^n}=\{z\in\mathbb{C}\mid z^{p^n}=1\},
\qquad U_{p^n}\leq(\mathbb{C}^{*},\cdot).
\]
\noindent Also $\operatorname{ord}(U_{p^n})=p^n$.

\[
\bigcup_{n\geq 1}U_{p^n}\ \leq\ (\mathbb{C}^{*},\cdot).
\]
\noindent This is an infinite $p$-group.

\bigskip
\noindent\underline{Def.}: Let $G$ be a group and $H\leq G$.
$H$ is a \underline{$p$-subgroup} of $G$ if and only if
$(H,\cdot)$ is a $p$-group.

\bigskip
\noindent\fbox{\textbf{Sylow's First Theorem}}.
Let $G$ be a finite group with $\operatorname{ord}(G)=p^t m$,
$t\geq 1$ and $p\nmid m$.

\medskip
\noindent Then for every $i=\overline{1,t}$ there exists $H\leq G$
with $|H|=p^i$.

\medskip
\noindent Moreover, if $1\leq i<t$, there exists $H\leq G$,
$|H|=p^i$, which is normal in a subgroup of $G$
of order $p^{i+1}$.

\bigskip
\noindent\underline{Proof}: For $i=1$, O.K. by Cauchy's theorem.

\medskip
\noindent Under the same hypotheses, a $p$-subgroup of $G$
of order $p^t$ is called a \underline{Sylow $p$-subgroup} of $G$.
\quad (Notation: $A=\operatorname{Syl}_p(G)$.)

\bigskip
\noindent\underline{Remark}: On $A=\operatorname{Syl}_p(G)$,
$G\times A\longrightarrow A$ acts by conjugation:
\[
P\in A,\ a\in G,\qquad
(a,P)\longmapsto aPa^{-1}
\quad\text{(the conjugate of $P$ by $a$)}.
\]
.
% No preamble and no forced page break.

\subsection*{Lecture 10 - Page 6}

\noindent\underline{Def.}: A group $G$ is called a \textbf{$p$-group},
$p$ prime, if:
\[
(\forall)\,a\in G,\qquad o(a)=p^t,\quad t\geq 0.
\]
% In the manuscript, the term "p-group" in the opening sentence is crossed out,
% although it is the term defined by the condition that follows.

\medskip
\noindent\underline{Remark}: If $G$ is a finite group, then
\[
\boxed{G\text{ is a }p\text{-group}\quad\Longleftrightarrow\quad
\operatorname{ord}(G)=p^n}.
\]

\noindent\underline{Proof}: ``$\Leftarrow$'' evident.

\noindent ``$\Rightarrow$'': By contradiction, suppose
$\operatorname{ord}(G)\neq p^n$.
Then there is a prime $q\neq p$ such that $q\mid\operatorname{ord}(G)$.

\noindent By Cauchy's theorem, $(\exists)\,a\in G$ with $o(a)=q$;
since $G$ is a $p$-group, this gives $p=q$, a contradiction.

\bigskip
\noindent-- Let
\[
U_{p^n}=\{z\in\mathbb{C}\mid z^{p^n}=1\},
\qquad U_{p^n}\leq(\mathbb{C}^{*},\cdot).
\]
\noindent Also $\operatorname{ord}(U_{p^n})=p^n$.

\[
\bigcup_{n\geq 1}U_{p^n}\ \leq\ (\mathbb{C}^{*},\cdot).
\]
\noindent This is an infinite $p$-group.

\bigskip
\noindent\underline{Def.}: Let $G$ be a group and $H\leq G$.
$H$ is a \underline{$p$-subgroup} of $G$ if and only if
$(H,\cdot)$ is a $p$-group.

\bigskip
\noindent\fbox{\textbf{Sylow's First Theorem}}.
Let $G$ be a finite group with $\operatorname{ord}(G)=p^t m$,
$t\geq 1$ and $p\nmid m$.

\medskip
\noindent Then for every $i=\overline{1,t}$ there exists $H\leq G$
with $|H|=p^i$.

\medskip
\noindent Moreover, if $1\leq i<t$, there exists $H\leq G$,
$|H|=p^i$, which is normal in a subgroup of $G$
of order $p^{i+1}$.

\bigskip
\noindent\underline{Proof}: For $i=1$, O.K. by Cauchy's theorem.

\medskip
\noindent Under the same hypotheses, a $p$-subgroup of $G$
of order $p^t$ is called a \underline{Sylow $p$-subgroup} of $G$.
\quad (Notation: $A=\operatorname{Syl}_p(G)$.)

\bigskip
\noindent\underline{Remark}: On $A=\operatorname{Syl}_p(G)$,
$G\times A\longrightarrow A$ acts by conjugation:
\[
P\in A,\ a\in G,\qquad
(a,P)\longmapsto aPa^{-1}
\quad\text{(the conjugate of $P$ by $a$)}.
\]
